\documentclass[12pt]{article}
\usepackage[utf8]{inputenc}
\usepackage[T1]{fontenc}
\usepackage[margin=1in]{geometry}
\usepackage{setspace}
\usepackage{xcolor}
\usepackage[colorlinks=true, urlcolor=blue!60!black, linkcolor=blue!60!black]{hyperref}
\setlength{\parindent}{2em}
\setlength{\parskip}{1.25em}
\renewcommand{\baselinestretch}{1.0}
\begin{document}
\begin{center}
{\LARGE \textit{Agency for International Development v.\ Alliance for Open Society International, Inc.}, 570 U.S.\ 205 (2013)\par}
\vspace{0.8em}
{\large Case Brief\par}
\end{center}
\vspace{1em}

\section{Memory Jogger}
Under the Spending Clause and the First Amendment's unconstitutional conditions doctrine, Congress may specify the activities it wishes to subsidize within a federal program, but it cannot compel domestic grant recipients to adopt and espouse the Government's viewpoint as their own belief outside the contours of the funded program (\href{https://supreme.justia.com/cases/federal/us/570/205/}{570 U.S. at 214--15, 221}).

\section{Detailed Case Facts}
In 2003, responding to a global HIV/AIDS pandemic that had infected more than 65 million people and killed over 25 million worldwide (including more than 19 million in sub-Saharan Africa), Congress enacted the United States Leadership Against HIV/AIDS, Tuberculosis, and Malaria Act of 2003 (the ``Leadership Act''), Pub. L. No. 108-25, 117 Stat. 711, codified at 22 U.S.C. \S~7601 et seq. (commonly associated with PEPFAR) (\href{https://supreme.justia.com/cases/federal/us/570/205/}{570 U.S. at 208--09}). Congress authorized billions of dollars in funding for nongovernmental organizations (NGOs) to combat HIV/AIDS abroad, declaring that the sex industry, trafficking, and sexual violence fueled the epidemic and that ``it should be the policy of the United States to eradicate'' prostitution and other sexual victimization (22 U.S.C. \S~7601(23)) (\href{https://supreme.justia.com/cases/federal/us/570/205/}{570 U.S. at 209}).

Congress attached two related conditions to Leadership Act grants:
\begin{itemize}
  \item \textbf{Use-of-Funds Restriction (22 U.S.C. \S~7631(e)):} No funds made available under the Act ``may be used to promote or advocate the legalization or practice of prostitution or sex trafficking.'' This condition was \emph{not} challenged (\href{https://supreme.justia.com/cases/federal/us/570/205/}{570 U.S. at 208, 217}).
  \item \textbf{Policy Requirement (22 U.S.C. \S~7631(f)):} No funds may be used to ``provide assistance to any group or organization that does not have a policy explicitly opposing prostitution and sex trafficking,'' exempting certain international bodies such as the World Health Organization, the Global Fund, and United Nations agencies (\href{https://supreme.justia.com/cases/federal/us/570/205/}{570 U.S. at 209--10}).
\end{itemize}
To enforce \S~7631(f), the Department of Health and Human Services (HHS) and the United States Agency for International Development (USAID) required every recipient to agree in its award documents that it is opposed to ``prostitution and sex trafficking because of the psychological and physical risks they pose for women, men, and children'' (45 C.F.R. \S~89.1(b); USAID Acquisition \& Assistance Policy Directive 12-04) (\href{https://supreme.justia.com/cases/federal/us/570/205/}{570 U.S. at 210}).

Respondents---Alliance for Open Society International, Inc. (AOSI), Pathfinder International, Global Health Council, and InterAction---were domestic U.S. NGOs receiving both substantial private donations and billions in federal aid to run overseas HIV/AIDS programs (such as reducing injection drug use in Uzbekistan, Tajikistan, and Kyrgyzstan; preventing mother-to-child transmission in Kenya; and promoting safer sex in India) (\href{https://supreme.justia.com/cases/federal/us/570/205/}{570 U.S. at 210--11}; \href{https://www.oyez.org/cases/2012/12-10}{Oyez}). Respondents did not promote prostitution, but wished to remain neutral because explicitly condemning prostitution risked alienating host governments, undermining trust with sex workers needed for effective HIV prevention, and censoring their privately funded public-health advocacy (\href{https://supreme.justia.com/cases/federal/us/570/205/}{570 U.S. at 211}; \href{https://www.law.cornell.edu/supct/cert/12-10}{Cornell LII}).

\textbf{Constitutional Provisions:} The case pitted Congress's Spending Clause authority to ``provide for the common Defence and general Welfare of the United States'' (U.S. Const. art. I, \S~8, cl. 1) against the First Amendment's Free Speech Clause (``Congress shall make no law \ldots abridging the freedom of speech'') and the doctrine of unconstitutional conditions (\href{https://www.law.cornell.edu/constitution-conan/amendment-1/conditions-exceeding-the-scope-of-the-program}{Constitution Annotated}).

\section{Procedural History}
In 2005, AOSI and Pathfinder International sued USAID, HHS, and federal officials in the U.S. District Court for the Southern District of New York, seeking declaratory relief and a preliminary injunction against enforcement of \S~7631(f) (\href{https://supreme.justia.com/cases/federal/us/570/205/}{570 U.S. at 211}). The District Court granted a preliminary injunction in 2006, 430 F. Supp. 2d 222 (S.D.N.Y. 2006) (\href{https://law.justia.com/cases/federal/appellate-courts/ca2/15-974/15-974-2018-12-20.html}{Justia 2d Cir. History}).

While the Government's appeal was pending, HHS and USAID issued ``affiliate guidelines'' (45 C.F.R. \S~89.3; AAPD 12-04) permitting grant recipients to work with affiliated entities that engaged in speech inconsistent with an anti-prostitution stance, provided the recipient maintained ``objective integrity and independence'' measured by five factors: legal separation, separate personnel, separate accounting, separate facilities, and distinct signage/identification (\href{https://supreme.justia.com/cases/federal/us/570/205/}{570 U.S. at 211}). The Second Circuit summarily remanded; the District Court reissued the preliminary injunction, and the Second Circuit affirmed (2--1), 651 F.3d 218 (2d Cir. 2011), holding that the Policy Requirement unconstitutionally compelled affirmative viewpoint speech, over Judge Straub's dissent (\href{https://supreme.justia.com/cases/federal/us/570/205/}{570 U.S. at 212}). Because the D.C. Circuit had previously upheld \S~7631(f) in \textit{DKT International, Inc. v. USAID}, 477 F.3d 758 (D.C. Cir. 2007), the Supreme Court granted certiorari on January 11, 2013, to resolve the circuit split (\href{https://www.everycrsreport.com/reports/R43027.html}{CRS Report R43027}). Argued April 22, 2013, and decided June 20, 2013, the Court affirmed (\href{https://www.oyez.org/cases/2012/12-10}{Oyez}).

\textbf{Subsequent History:} In \href{https://supreme.justia.com/cases/federal/us/591/19-177/}{\textit{USAID v. Alliance for Open Society International, Inc.} (\textit{AOSI II}), 591 U.S. 430 (2020)}, the Supreme Court held (5--3) that the Policy Requirement remains constitutional as applied to plaintiffs' separately incorporated \emph{foreign} affiliates operating abroad, because foreign entities outside U.S. territory possess no First Amendment rights (\href{https://www.oyez.org/cases/2019/19-177}{Oyez}).

\section{Judicial Votes}
\begin{itemize}
  \item \textbf{Judgment (6--2):} Affirmed.
  \item \textbf{Majority Opinion:} Chief Justice Roberts, joined by Justices Kennedy, Ginsburg, Breyer, Alito, and Sotomayor (\href{https://supreme.justia.com/cases/federal/us/570/205/}{570 U.S. at 207}).
  \item \textbf{Concurrences:} None.
  \item \textbf{Dissenting Opinion:} Justice Scalia, joined by Justice Thomas (\href{https://supreme.justia.com/cases/federal/us/570/205/}{570 U.S. at 221}).
  \item \textbf{Recusal:} Justice Kagan took no part in the consideration or decision of the case (\href{https://supreme.justia.com/cases/federal/us/570/205/}{570 U.S. at 221}).
\end{itemize}

\section{Holding}
\textbf{Judgment:} The judgment of the Second Circuit affirming the preliminary injunction against enforcing 22 U.S.C. \S~7631(f) as to domestic U.S. organizations is affirmed (\href{https://supreme.justia.com/cases/federal/us/570/205/}{570 U.S. at 221}).

\textbf{Rule Announced:} The First Amendment's prohibition on compelled speech bars Congress from conditioning federal funding on a domestic recipient's affirmative adoption and espousal of a government viewpoint that, by its nature, cannot be confined within the scope of the federally funded program (\href{https://supreme.justia.com/cases/federal/us/570/205/}{570 U.S. at 214--15, 221}). Under the unconstitutional conditions doctrine, courts distinguish between permissible conditions that \emph{define the limits of the government spending program} (specifying the activities Congress chooses to subsidize) and impermissible conditions that \emph{leverage funding to regulate speech outside the contours of the program itself} (\href{https://supreme.justia.com/cases/federal/us/570/205/}{570 U.S. at 214--15}).

\textbf{Holding vs. Dicta:} The holding invalidates \S~7631(f)'s Policy Requirement as applied to domestic NGOs and rejects the Government's affiliate-guidelines cure and fungibility defense on this record. Notably, the Court's acknowledgment that the program-versus-outside-the-program line is ``hardly clear'' and that ``Definition more precise must abide the wisdom of the future'' (quoting \textit{Steward Machine Co. v. Davis}, 301 U.S. 548, 591 (1937)) is framing commentary rather than a bright-line rule (\href{https://supreme.justia.com/cases/federal/us/570/205/}{570 U.S. at 215, 217}). Additionally, the opinion did not decide whether \S~7631(f) could be applied to foreign NGOs---the precise issue later resolved in the Government's favor in \href{https://supreme.justia.com/cases/federal/us/591/19-177/}{\textit{AOSI II} (2020)}.

\section{Analysis of Opinions}

\subsection{Majority Opinion (Roberts, C.J.)}
Chief Justice Roberts opened with bedrock First Amendment compelled-speech doctrine: ``freedom of speech prohibits the government from telling people what they must say'' (\textit{Rumsfeld v. FAIR}, 547 U.S. 47, 61 (2006); \textit{West Virginia State Bd. of Ed. v. Barnette}, 319 U.S. 624 (1943); \textit{Wooley v. Maynard}, 430 U.S. 705 (1977)) (\href{https://supreme.justia.com/cases/federal/us/570/205/}{570 U.S. at 213}). As a direct regulation, the Policy Requirement would plainly violate the First Amendment. The sole question was whether Congress could achieve the same result indirectly under the Spending Clause.

Roberts synthesized the Court's unconstitutional-conditions precedent into the ``inside the program'' versus ``outside the program'' distinction (\href{https://supreme.justia.com/cases/federal/us/570/205/}{570 U.S. at 214--17}):
\begin{itemize}
  \item In \textit{Regan v. Taxation With Representation of Washington}, 461 U.S. 540 (1983), barring \S~501(c)(3) charities from substantial lobbying was valid because Congress merely declined to subsidize lobbying, and the group could freely lobby using private funds through a dual \S~501(c)(4) affiliate.
  \item In \textit{FCC v. League of Women Voters of California}, 468 U.S. 364 (1984), barring federally funded public broadcasters from \emph{all} editorializing---even with private funds, when federal money was just 1\% of the budget---was struck down because it leveraged funding to control speech outside the program.
  \item In \textit{Rust v. Sullivan}, 500 U.S. 173 (1991), barring Title X family-planning projects from abortion counseling was upheld because the regulations distinguished between the \emph{Title X project} and the \emph{Title X grantee}, leaving the grantee unfettered to advocate abortion in separate, privately funded programs.
\end{itemize}
Applying this framework, Roberts found the Policy Requirement squarely on the unconstitutional side of the line. Because unchallenged \S~7631(e) already ensured no federal funds promoted prostitution, \S~7631(f) necessarily went further: it demanded that recipients adopt the Government's orthodoxy as their own institutional belief, thereby ``defining the recipient'' rather than ``defining the limits of the federally funded program'' (\href{https://supreme.justia.com/cases/federal/us/570/205/}{570 U.S. at 217--18}). An NGO cannot avow opposition to prostitution while spending federal funds and claim neutrality or dissent on its own dime.

Roberts rejected three government defenses: (1) the \textbf{affiliate guidelines} failed because an independent affiliate cannot let the \emph{recipient} express its own beliefs, while a closely identified affiliate forces the recipient into ``evident hypocrisy'' (\href{https://supreme.justia.com/cases/federal/us/570/205/}{570 U.S. at 219}); (2) the \textbf{money-is-fungible} theory (invoking \textit{Holder v. Humanitarian Law Project}, 561 U.S. 1 (2010)) lacked any record evidence that federal grants merely supplanted private funds and would have overturned \textit{League of Women Voters} (\href{https://supreme.justia.com/cases/federal/us/570/205/}{570 U.S. at 219--20}); and (3) preventing ``message confusion'' could not justify forcing grantees ``to pledge allegiance to the Government's policy of eradicating prostitution,'' violating \textit{Barnette}'s fixed-star principle (\href{https://supreme.justia.com/cases/federal/us/570/205/}{570 U.S. at 220--21}).

\subsection{Dissenting Opinion (Scalia, J., joined by Thomas, J.)}
Justice Scalia characterized \S~7631(f) as a commonsense selection criterion enabling the Government to enlist ideological allies to implement its chosen HIV/AIDS strategy (\href{https://supreme.justia.com/cases/federal/us/570/205/}{570 U.S. at 221--22}). ``The First Amendment does not mandate a viewpoint-neutral government,'' Scalia argued; just as the United States need not hire Hamas to distribute foreign food aid or the American Gourmet Society to run a healthy-eating campaign, Congress may choose grantees who share its anti-prostitution mission (\href{https://supreme.justia.com/cases/federal/us/570/205/}{570 U.S. at 221--22}).

Scalia attacked the majority's inside/outside-the-program test as unmoored from the First Amendment. Unlike \textit{Rust}, where Title X was neutral on abortion outside the clinic, eradicating prostitution was a core objective of the Leadership Act, so outside promotion undermined the program itself---especially because ``[m]oney is fungible'' (\href{https://supreme.justia.com/cases/federal/us/570/205/}{570 U.S. at 223--24}). Crucially, Scalia argued there was no \emph{coercion}: ``the Government is not forcing anyone to say anything,'' because any NGO may simply decline the grant and remain as unconstrained as before (\href{https://supreme.justia.com/cases/federal/us/570/205/}{570 U.S. at 225--27}). Invoking the Article VI Oath Clause (U.S. Const. art. VI, cl. 3), Scalia contended that the Framers themselves approved requiring relevant affirmative ideological commitments as a prerequisite to assisting in the Government's work, and warned that the majority's rule would invite endless lawsuits whenever the Government selects grantees on relevant ideological grounds (\href{https://supreme.justia.com/cases/federal/us/570/205/}{570 U.S. at 226--27}).

\section{Examples: Future Applications}
\begin{enumerate}
  \item \textbf{Same-Side (Unconstitutional Condition):} Congress appropriates federal grants for domestic opioid-overdose prevention clinics but bars funding to any organization that does not adopt an official corporate policy opposing the decriminalization of all controlled substances. Even if Congress also bans using grant money to lobby for decriminalization, requiring the clinic itself to adopt an anti-decriminalization creed reaches beyond the funded program to compel institutional belief outside it, violating \textit{USAID v. AOSI}.
  \item \textbf{Same-Side (Unconstitutional Condition):} A state conditions eligibility for general public-university research grants on each university's formal public pledge endorsing the state legislature's foreign-policy stance on an international boycott. Because the pledge dictates institutional orthodoxy unrelated to defining the specific research project being subsidized, it unlawfully leverages funding to compel speech outside the program's contours (\href{https://houstonlawreview.org/article/175172-coercion-conditions-and-campus-governance-columbia-university-and-the-constitutional-boundaries-of-federal-funding}{Houston L. Rev. (2026)}).
  \item \textbf{Opposite-Side (Permissible Program Definition):} Congress funds a federal ``Clean Energy Youth Corps'' grant program and requires that all federally funded curriculum materials and classroom presentations teach that solar and wind power reduce carbon emissions, while leaving grantee nonprofits completely free to publish skeptical articles on their own websites using private funds. Because the restriction defines only the message delivered \emph{within} the subsidized program, it is constitutional under \textit{Rust} and \textit{USAID v. AOSI}.
  \item \textbf{Opposite-Side (Permissible Foreign Application):} Congress conditions foreign-aid grants awarded directly to independently incorporated European or African NGOs operating exclusively overseas on adopting a policy opposing human trafficking and commercial sex work. Under \href{https://supreme.justia.com/cases/federal/us/591/19-177/}{\textit{AOSI II}, 591 U.S. 430 (2020)}, because foreign organizations outside U.S. territory possess no First Amendment rights, this condition is valid even though it mirrors \S~7631(f).
  \item \textbf{Fence-Sitter (Borderline Case):} A federal agency awards competitive grants for an ``Anti-Smoking Media Campaign'' and scores applications partly on whether the applicant organization has a demonstrated organizational commitment to tobacco cessation, disqualifying a public-health NGO that accepts tobacco-industry sponsorships in its other divisions. This sits squarely on the razor's edge highlighted by Scalia's dissent: is the agency permissibly selecting a credible messenger to define the program, or impermissibly penalizing off-program speech and association (\href{https://supreme.justia.com/cases/federal/us/570/205/}{570 U.S. at 226})?
\end{enumerate}

\section{Critique}
\textit{USAID v. AOSI} was widely welcomed for reinforcing \textit{Barnette}'s anti-orthodoxy principle in the modern grant-making state, yet scholars across the ideological spectrum criticize its doctrinal indeterminacy (\href{https://harvardlawreview.org/print/vol-127/first-amendment-ae-freedom-of-speech-ae-unconstitutional-conditions-ae-agency-for-international-development-v-alliance-for-open-society-international-inc/}{127 Harv. L. Rev. 218 (2013)}).

\textbf{Logical Weaknesses and Manipulability:} Chief Justice Roberts candidly conceded that the distinction between conditions ``defining the program'' and those ``outside the program'' is ``hardly clear,'' because Congress can often redefine a program broadly enough to subsume the condition (\href{https://supreme.justia.com/cases/federal/us/570/205/}{570 U.S. at 215}). Moreover, the Court's opinion never squarely answers Justice Scalia's sharpest practical challenge: if the Government cannot require an applicant to \emph{adopt} a policy as a condition of funding, may it still \emph{prefer} or \emph{select} applicants who already hold that view during competitive grantmaking? If so, the Court's rule may elevate form (explicit pledge requirements) over substance (covert ideological selection).

\textbf{Originalist and Textualist Perspectives:} Textualists and originalists split over \textit{USAID v. AOSI}. Scalia and Thomas's dissent reflects a strict textual reading of the First Amendment's word ``abridging'': refusing to subsidize a private actor who remains legally free to decline federal funds is not ``coercion'' or ``abridgment'' of speech, and the Article VI Oath Clause confirms that ideological alignment with government missions was historically acceptable (\href{https://firstamendment.mtsu.edu/article/agency-for-international-development-v-alliance-for-open-society/}{First Amendment Encyclopedia}). Conversely, libertarian originalists and Spending Clause skeptics applaud Roberts's result for checking Congress's vast modern spending power, arguing that without a robust unconstitutional conditions doctrine, the federal government could buy up constitutional rights that the enumerated-powers framework forbids it from regulating directly.

\textbf{Progressive and Pragmatic Perspectives:} Progressive scholars generally praised \textit{AOSI I} for protecting public-health NGOs' autonomy and preventing ideological distortion of evidence-based HIV prevention (\href{https://www.law.cornell.edu/supct/cert/12-10}{Cornell LII}). However, many progressive critics argue that Roberts's formalistic reliance on corporate separation left domestic NGOs vulnerable to the later retreat in \href{https://supreme.justia.com/cases/federal/us/591/19-177/}{\textit{AOSI II} (2020)}, where the Court allowed the Government to enforce the exact same pledge against closely branded foreign affiliates---effectively forcing U.S. NGOs into the very ``evident hypocrisy'' Roberts had condemned in 2013 (\href{https://www.americanbar.org/groups/public_contract_law/resources/journal/2025-spring/preserving-government-contractor-speech-rights-overlooked-implications/}{ABA Public Contract L.J. (2025)}). Progressive theorists also worry that an expansive compelled-speech doctrine in funding cases can be weaponized to challenge federal nondiscrimination and labor conditions attached to public grants and contracts.

\section{Key Quotations}
\begin{quote}
``In the present context, the relevant distinction that has emerged from our cases is between conditions that define the limits of the government spending program---those that specify the activities Congress wants to subsidize---and conditions that seek to leverage funding to regulate speech outside the contours of the program itself.''
\end{quote}
\noindent --- Chief Justice Roberts, Opinion of the Court (\href{https://supreme.justia.com/cases/federal/us/570/205/}{570 U.S. at 214--15}).

\begin{quote}
``By demanding that funding recipients adopt---as their own---the Government's view on an issue of public concern, the condition by its very nature affects `protected conduct outside the scope of the federally funded program.' \ldots By requiring recipients to profess a specific belief, the Policy Requirement goes beyond defining the limits of the federally funded program to defining the recipient.''
\end{quote}
\noindent --- Chief Justice Roberts, Opinion of the Court (\href{https://supreme.justia.com/cases/federal/us/570/205/}{570 U.S. at 218}).

\begin{quote}
``Affiliates cannot serve that purpose when the condition is that a funding recipient espouse a specific belief as its own. If the affiliate is distinct from the recipient, the arrangement does not afford a means for the recipient to express its beliefs. If the affiliate is more clearly identified with the recipient, the recipient can express those beliefs only at the price of evident hypocrisy.''
\end{quote}
\noindent --- Chief Justice Roberts, Opinion of the Court (\href{https://supreme.justia.com/cases/federal/us/570/205/}{570 U.S. at 219}).

\begin{quote}
``But the Policy Requirement goes beyond preventing recipients from using private funds in a way that would undermine the federal program. It requires them to pledge allegiance to the Government's policy of eradicating prostitution. As to that, we cannot improve upon what Justice Jackson wrote for the Court 70 years ago: `If there is any fixed star in our constitutional constellation, it is that no official, high or petty, can prescribe what shall be orthodox in politics, nationalism, religion, or other matters of opinion or force citizens to confess by word or act their faith therein.' ''
\end{quote}
\noindent --- Chief Justice Roberts, Opinion of the Court (\href{https://supreme.justia.com/cases/federal/us/570/205/}{570 U.S. at 220--21}).

\begin{quote}
``The First Amendment does not mandate a viewpoint-neutral government. Government must choose between rival ideas and adopt some as its own: competition over cartels, solar energy over coal, weapon development over disarmament, and so forth. Moreover, the government may enlist the assistance of those who believe in its ideas to carry them to fruition; and it need not enlist for that purpose those who oppose or do not support the ideas.''
\end{quote}
\noindent --- Justice Scalia, dissenting (\href{https://supreme.justia.com/cases/federal/us/570/205/}{570 U.S. at 221--22}).


\end{document}
